\documentclass[11pt]{article}
\usepackage[margin=25mm]{geometry}
\usepackage{fontspec}\setmainfont{Latin Modern Roman}
\usepackage{amsmath,amssymb,mathtools}
\usepackage{unicode-math}\setmathfont{Latin Modern Math}
\usepackage{xurl,hyperref}
\hypersetup{colorlinks=true,urlcolor=blue}
\setcounter{secnumdepth}{0}
\setlength{\parindent}{0pt}\setlength{\parskip}{5pt}
\setlength{\emergencystretch}{3em}
\newcommand{\sref}[2]{\hyperlink{source:#1:#2}{[S#2]}}
\newcommand{\cataloguescope}[1]{\paragraph{Recorded target.}#1}
\begin{document}
\section{Kourovka Notebook Problem 21.134}\label{statement}

\subsection{Definitions and mathematical statement}

\paragraph{Shared notation.}$\mathbb N=\{1,2,\ldots\}$, and $\mathbb Z,\mathbb Q,\mathbb R,\mathbb C$ denote the integers, rationals, reals and complex numbers. Any other conventions are specified locally.


For a finite group $G$, put $t_G(n)=|\{x\in G:x^n=1\}|$ for $n\geq1$. The solvable radical $R(G)$ is the largest solvable normal subgroup. A finite group $H$ is almost simple if $S\leq H\leq\operatorname{Aut}(S)$ for some nonabelian finite simple group $S$.
\paragraph{Statement recorded in the supplied source.}
Suppose finite groups $G,H$ satisfy $t_G(n)=t_H(n)$ for every $n\geq1$. (a) If $R(H)=1$, must $R(G)=1$? (b) If $H$ is almost simple, must $G\cong H$? \sref{statement}{1}\sref{statement}{2}
\subsection{Short English statement}
Ask whether the numbers of solutions to every power equation detect a trivial solvable radical and uniquely identify an almost simple group.
\subsection{Sources}
\begin{description}
\item[\hypertarget{source:statement:1}{S1}] ULAM.AI and the UnsolvedMath Contributors, \emph{UnsolvedMath: A Curated Collection of Open Mathematics Problems}, catalogue entry. CC BY 4.0. \url{https://huggingface.co/datasets/ulamai/UnsolvedMath}.
\item[\hypertarget{source:statement:2}{S2}] E. I. Khukhro and V. D. Mazurov (editors), The Kourovka Notebook, Unsolved Problems in Group Theory, No.21, version48 (6 October2026), Problem 21.134 and its remarks. \url{https://arxiv.org/abs/1401.0300v48}.
\end{description}
\label{end:statement}
\end{document}
